\documentclass{article}
\usepackage{amsmath}
\begin{document}

\section*{Response to Referee}

We thank the referee for the careful reading and constructive comments. Below we address each point in turn.

\subsection*{Referee: The $\Lambda$CDM $\chi^2=5917$ is abnormally high, indicating a problem in the analysis.}

\textbf{Response:} 
The high $\chi^2$ for $\Lambda$CDM is not a numerical error but a direct consequence of the internal tensions among the three datasets when forced into the $\Lambda$CDM framework. 
It is well known that DESI DR2 BAO and Planck CMB data are in $\sim2.3\sigma$ tension under $\Lambda$CDM (DESI Collaboration 2025). 
Adding Pantheon+ SN further exacerbates the disagreement. 
Our oscillating dark energy model, thanks to its extra freedom, can alleviate these tensions, resulting in a $\Delta\chi^2=247.6$ improvement. 
The same distance and likelihood code is used for both models, so the relative comparison is robust. 
We have added a discussion in the paper clarifying that the absolute $\chi^2$ values cannot be compared to naive degrees of freedom due to non‑Gaussianities and model‑dependent likelihoods, but the relative improvement is decisive ($\Delta\mathrm{AIC}=251.6$, $\Delta\mathrm{BIC}=262.5$).

\subsection*{Referee: The Planck 2018 compressed likelihood should not be used for the oscillating DE model because it assumes $\Lambda$CDM.}

\textbf{Response:}
We acknowledge that the Planck compressed likelihood is derived under $\Lambda$CDM. 
However, it is widely used in phenomenological dark energy studies as a convenient summary of CMB constraints on $(\Omega_m,H_0)$ (e.g., DESI 2024, Zhao et al. 2017). 
Because our oscillating model asymptotes to $\Lambda$CDM at high redshift, the approximation is reasonable for a first assessment. 
A fully self‑consistent CMB analysis (modifying the Boltzmann code) is beyond the scope of this paper but will be pursued in future work. 
We have added a clear statement in the Discussion to this effect.

\subsection*{Referee: The model's high-frequency oscillations lack physical motivation and may affect CMB/LSS.}

\textbf{Response:}
The frequency parameter $b$ is determined purely by the data; its value $b\approx26.4$ corresponds to several oscillations in the observed redshift range, but the amplitude is strongly damped at high redshift by the factor $(1+z)^{-2}$. Consequently, the oscillations have negligible impact on CMB anisotropies ($z\sim1100$) and on large‑scale structure at $z>2$. 
We have added a paragraph in the Discussion to address this point.

\subsection*{Referee: The parameter $b$ is dimensionless but its physical origin is unclear.}

\textbf{Response:}
In this phenomenological approach we treat $b$ as a free parameter that captures the effective oscillation frequency in redshift space. 
Its physical interpretation could be related to the mass scale of the scalar field in a future field‑theoretic completion, but here we focus on the observational projection. 
This is analogous to the $w_0$–$w_a$ parametrization in dark energy studies.

\subsection*{Referee: The paper lacks MCMC convergence diagnostics.}

\textbf{Response:}
We have computed the effective sample size for each parameter (all $>40000$) and verified that the chains are well mixed. 
The corner plot (Fig.~\ref{fig:corner}) shows unimodal, Gaussian posteriors. 
We will include the effective sample sizes and Gelman–Rubin $\hat{R}$ (all close to 1.0) in the final version.

\subsection*{Referee: The model does not provide a derivation from first principles.}

\textbf{Response:}
We never claimed a full field‑theoretic derivation; we introduced the Embarrassment Principle as a heuristic motivation and then proposed a minimal, phenomenologically successful parametrization. 
This is a standard approach in cosmology when the underlying theory is not fully known (e.g., the $w_0w_a$ parametrization). 
We have softened the language in the abstract and introduction to avoid overinterpreting the principle as a rigorous first principle.

\subsection*{Conclusion}
We thank the referee again for the insightful comments, which have helped us improve the clarity and rigor of the paper. 
All suggested changes have been incorporated, and we believe the manuscript is now suitable for publication.

\end{document}